% !TEX root = ../proyect.tex

\chapter{Resultados y discusión}\label{cap:resultados}

\section{Resultados de la comprobación sintética}\label{sec:resultados-comprobacion-sintetica}

La comprobación sintética se utiliza como control metodológico del procedimiento. A diferencia de la serie financiera, el mapa logístico proporciona una dinámica conocida, determinista y no lineal. Por tanto, permite evaluar si las herramientas de reconstrucción, cuantificación y predicción local responden de forma coherente cuando la estructura subyacente no es una hipótesis abierta, sino una propiedad del sistema generado.

El sistema utilizado es el mapa logístico en el caso \(r=4\),
\begin{equation}
    x_{t+1}=4x_t(1-x_t),
\end{equation}
con \(x_0=0{,}123456789\). Se generan 12.000 iteraciones y se descartan las primeras 1.000 como transitorio, quedando 11.000 observaciones útiles. Además de la serie limpia, se construyen dos versiones con ruido observacional: una con ruido pequeño y otra con ruido moderado. El ruido no modifica la regla determinista subyacente, pero degrada la observación disponible para el procedimiento.

\begin{figure}[htbp]
    \centering
    \includegraphics[width=0.78\textwidth]{figures/phase13_transition_map_clean}
    \caption{Mapa de transición de la serie logística limpia. La nube de pares \((x_t,x_{t+1})\) sigue la parábola teórica \(4x(1-x)\), mientras que la recta ajustada sobre entrenamiento ilustra la limitación de una aproximación lineal global.}
    \label{fig:phase13-transition-map-clean}
\end{figure}

La Figura~\ref{fig:phase13-transition-map-clean} muestra el elemento central de la comprobación: la transición real es cuadrática, no lineal y acotada en \([0,1]\). La curva teórica se representa solo como referencia visual; no se usa para entrenar los modelos. La comparación con la recta ajustada sobre entrenamiento permite anticipar por qué una referencia AR lineal puede servir como contraste, pero no como descripción estructural adecuada del sistema.

\subsection{Reconstrucción de la dinámica sintética}

La selección de parámetros de reconstrucción se realiza de forma análoga a la serie financiera. La información mutua media selecciona \(\tau=9\) para la serie limpia y para la serie con ruido pequeño, mientras que en la serie con ruido moderado el primer mínimo local aparece en \(\tau=5\). La dimensión práctica se fija mediante FNN, contrastada con Cao. En las tres variantes, Cao tiende hacia dimensiones más altas, por lo que la dimensión seleccionada debe interpretarse como operativa, no como estimación exacta de una dimensión intrínseca.

\begin{table}[htbp]
\centering
\small
\begin{tabular}{lrrrr}
\hline
\textbf{Serie} & \(\boldsymbol{\tau}\) & \(\boldsymbol{m_{\mathrm{FNN}}}\) & \(\boldsymbol{m_{\mathrm{Cao}}}\) & \(\boldsymbol{m}\) seleccionado \\
\hline
Limpia & 9 & 5 & 10 & 5 \\
Ruido pequeño & 9 & 6 & 10 & 6 \\
Ruido moderado & 5 & 5 & 10 & 5 \\
\hline
\end{tabular}
\caption{Parámetros de reconstrucción seleccionados en la comprobación sintética.}
\label{tab:phase13-embedding-params}
\end{table}

La discrepancia entre FNN y Cao no invalida la fase. El mapa logístico es un mapa discreto unidimensional no invertible, por lo que los criterios de reconstrucción por retardos no se comportan exactamente igual que en el caso ideal de un flujo suave invertible. Lo relevante para esta comprobación no es obtener una dimensión ``verdadera'', sino comprobar que la representación por retardos conserva estructura suficiente para cuantificación y predicción local.

\begin{figure}[htbp]
    \centering
    \includegraphics[width=0.72\textwidth]{figures/phase13_embedding_2d_clean}
    \caption{Reconstrucción bidimensional de la serie logística limpia con el retardo seleccionado. La nube reconstruida es acotada y no uniforme, lo que refleja una estructura geométrica más definida que la esperable en una serie puramente aleatoria.}
    \label{fig:phase13-embedding-clean}
\end{figure}

La Figura~\ref{fig:phase13-embedding-clean} muestra la reconstrucción bidimensional de la serie limpia. No aparece una curva simple, porque el retardo seleccionado separa coordenadas en una dinámica fuertemente mezclante. Aun así, la nube no es arbitraria: queda limitada por la propia dinámica del mapa y presenta concentraciones no uniformes. Al añadir ruido, según los informes de la fase, la forma global se mantiene en el caso de ruido pequeño y se difumina con mayor claridad en el caso de ruido moderado.

\subsection{Divergencia local de trayectorias}

La estimación tipo Rosenstein se utiliza como indicador cualitativo de divergencia local. Para \(r=4\), el exponente de Lyapunov teórico del mapa logístico es \(\ln(2)\approx 0{,}693\) por iteración. Las pendientes estimadas no deben compararse de forma directa con ese valor como si fueran una estimación exacta, ya que dependen del embedding, del bloque elegido y del intervalo de ajuste. La lectura relevante es la presencia o ausencia de una región inicial de crecimiento.

\begin{table}[htbp]
\centering
\small
\begin{tabular}{lrrrr}
\hline
\textbf{Serie} & \(\boldsymbol{\tau}\) & \(\boldsymbol{m}\) & \textbf{Pendiente} & \(\boldsymbol{R^2}\) ajuste \\
\hline
Limpia & 9 & 5 & 0{,}1852 & 0{,}6862 \\
Ruido pequeño & 9 & 6 & 0{,}1315 & 0{,}6331 \\
Ruido moderado & 5 & 5 & -0{,}0084 & 0{,}0014 \\
\hline
\end{tabular}
\caption{Resumen de la estimación de divergencia local mediante Rosenstein en las tres variantes del mapa logístico.}
\label{tab:phase13-lyapunov}
\end{table}

\begin{figure}[htbp]
    \centering
    \includegraphics[width=0.82\textwidth]{figures/phase13_lyapunov_rosenstein}
    \caption{Curvas de divergencia media tipo Rosenstein para la comprobación sintética. La serie limpia y la serie con ruido pequeño muestran crecimiento inicial positivo; la serie con ruido moderado no presenta una región lineal clara.}
    \label{fig:phase13-lyapunov}
\end{figure}

La Tabla~\ref{tab:phase13-lyapunov} y la Figura~\ref{fig:phase13-lyapunov} muestran una señal coherente con el propósito de la fase. En la serie limpia aparece una pendiente positiva; con ruido pequeño la pendiente sigue siendo positiva, aunque más baja; con ruido moderado desaparece una región de crecimiento lineal interpretable. Este comportamiento es razonable: el ruido observacional no elimina la dinámica subyacente, pero sí degrada la reconstrucción y dificulta la estimación de divergencia.

\subsection{Predicción local frente a referencias lineales}

La comprobación predictiva se plantea con una separación temporal 60--20--20 entre entrenamiento, validación y test. El horizonte es \(1\), porque el mapa define directamente \(x_{t+1}\) a partir de \(x_t\). En validación se selecciona \(k=5\) para las tres variantes. La referencia lineal se obtiene mediante selección AR por BIC entre \(p=0,\ldots,100\); en las tres series el criterio selecciona AR(0), equivalente a la media histórica.

\begin{figure}[htbp]
    \centering
    \includegraphics[width=0.82\textwidth]{figures/phase13_prediction_validation_k}
    \caption{Selección de \(k\) en validación para la predicción local del mapa logístico. El mínimo aparece en torno a \(k=5\) para las tres variantes.}
    \label{fig:phase13-validation-k}
\end{figure}

La Figura~\ref{fig:phase13-validation-k} indica que el error disminuye al pasar de un único vecino a una media local de varios vecinos, pero vuelve a aumentar cuando \(k\) crece demasiado. Este patrón es el esperado en un sistema sintético con estructura local marcada: pocos vecinos reducen la varianza respecto al vecino más próximo, pero demasiados vecinos mezclan estados menos parecidos y suavizan en exceso la regla de transición.

\begin{figure}[htbp]
    \centering
    \includegraphics[width=0.86\textwidth]{figures/phase13_prediction_real_vs_predicted_clean_compact}
    \caption{Predicción compacta en test para la serie logística limpia. El kNN local sigue la trayectoria real con mucha más precisión que persistencia y AR(0)/media.}
    \label{fig:phase13-prediction-clean}
\end{figure}

En la serie limpia, la Figura~\ref{fig:phase13-prediction-clean} muestra de forma directa la diferencia entre una referencia lineal global y una aproximación local. La persistencia falla porque \(x_{t+1}\) puede diferir notablemente de \(x_t\); el AR(0) queda cerca de la media; el kNN, en cambio, utiliza estados reconstruidos similares y aproxima de forma local la transición no lineal.

\begin{table}[htbp]
\centering
\small
\begin{tabular}{lrrrrr}
\hline
\textbf{Serie} & \textbf{AR/media RMSE} & \textbf{Persist. RMSE} & \textbf{NN RMSE} & \textbf{kNN RMSE} & \(\boldsymbol{R^2_{\mathrm{OOS}}}\) kNN \\
\hline
Limpia & 0{,}3565 & 0{,}4994 & 0{,}1272 & 0{,}0982 & 0{,}9242 \\
Ruido pequeño & 0{,}3564 & 0{,}4993 & 0{,}1794 & 0{,}1301 & 0{,}8667 \\
Ruido moderado & 0{,}3558 & 0{,}4985 & 0{,}1398 & 0{,}1075 & 0{,}9088 \\
\hline
\end{tabular}
\caption{Métricas de test en la comprobación sintética. En las tres variantes, el kNN medio seleccionado en validación obtiene el menor RMSE.}
\label{tab:phase13-prediction-metrics}
\end{table}

\begin{figure}[htbp]
    \centering
    \includegraphics[width=0.84\textwidth]{figures/phase13_prediction_metrics}
    \caption{Comparación de RMSE en test para la comprobación sintética. La media de vecinos mejora claramente a la media histórica, persistencia, AR(0) y vecino más próximo.}
    \label{fig:phase13-prediction-metrics}
\end{figure}

La Tabla~\ref{tab:phase13-prediction-metrics} y la Figura~\ref{fig:phase13-prediction-metrics} resumen el resultado más importante de la fase. En la serie limpia, el RMSE pasa de 0{,}3565 en AR(0)/media a 0{,}0982 con kNN. En la serie con ruido pequeño, el RMSE del kNN aumenta a 0{,}1301, pero sigue claramente por debajo de la referencia lineal. En la serie con ruido moderado, el RMSE queda en 0{,}1075. La degradación no es estrictamente monótona entre los dos niveles de ruido, pero en ambos casos el predictor local conserva una ventaja amplia frente a las referencias globales.

\subsection{Lectura metodológica de la comprobación}

La comprobación sintética cumple su función: en un sistema donde existe una regla no lineal conocida, el procedimiento produce una reconstrucción estructurada, detecta divergencia local en el caso limpio y permite una predicción local claramente superior a referencias lineales simples. Esto no implica que el mismo comportamiento deba aparecer en Bitcoin. La diferencia es precisamente el motivo de incluir el control: en el mapa logístico la señal no lineal domina la observación; en una serie financiera real, la volatilidad observada combina persistencia, ruido, heterocedasticidad, cambios de régimen y posibles componentes distribucionales.

Por tanto, el resultado de esta fase valida el funcionamiento del procedimiento en un entorno controlado, pero no demuestra nada por sí solo sobre Bitcoin. Su utilidad está en fijar un punto de comparación: cuando la dinámica no lineal es explícita, el método local en espacio reconstruido ofrece una ventaja clara; si en Bitcoin esa ventaja se debilita frente a modelos lineales o multiescala, la explicación debe buscarse en la naturaleza de la serie financiera y no necesariamente en un fallo de implementación.
\section{Resultados de la aplicación empírica a Bitcoin}\label{sec:resultados-aplicacion-bitcoin}
% TODO: redactar este apartado.

\section{Resultados de la caracterización inicial}\label{sec:resultados-caracterizacion}
% TODO: redactar este apartado.

\section{Resultados del filtrado lineal}\label{sec:resultados-filtrado}
% TODO: redactar este apartado.

\section{Indicios de dependencia no lineal}\label{sec:evidencia-no-lineal}
% TODO: redactar este apartado.

\section{Resultados de la reconstrucción del espacio de estados}\label{sec:resultados-reconstruccion}
% TODO: redactar este apartado.

\section{Resultados de cuantificación dinámica}\label{sec:resultados-cuantificacion}
% TODO: redactar este apartado.

\section{Resultados de predicción local}\label{sec:resultados-prediccion}
% TODO: redactar este apartado.

\section{Comparación con modelos de referencia}\label{sec:comparacion-referencia}
% TODO: redactar este apartado.

\section{Comparación entre configuraciones alternativas}\label{sec:comparacion-configuraciones}
% TODO: redactar este apartado.

\section{Comparación final con HAR-logRV}\label{sec:resultados-har}
La comparación final incorpora HAR-logRV como modelo práctico de memoria multiescala. En la muestra test comparable de 5000 puntos, HAR-logRV obtiene un RMSE de 0,8608, frente a 0,8641 de AR(49) y 0,8765 del kNN local con \(\tau=137\), \(m=5\) y \(k=200\). La mejora frente a AR(49) es pequeña, aproximadamente un 0,4\%, por lo que no debe presentarse como una ruptura fuerte respecto al benchmark lineal. Su interés principal está en que combina buen rendimiento, bajo coste computacional, interpretabilidad y facilidad de exportación al MVP.

Este resultado también matiza la lectura de las fases de reconstrucción: el kNN confirma cierta utilidad predictiva de la estructura local, pero en BTC la persistencia multiescala de la volatilidad queda mejor resumida por un modelo HAR simple. Por tanto, HAR-logRV se interpreta como cierre práctico del estudio, no como evidencia adicional de caos determinista.

\section{Discusión: indicios frente a demostración de caos determinista}\label{sec:discusion-caos}
% TODO: redactar este apartado.

\section{Interpretación global de los resultados}\label{sec:interpretacion-global}

kNN era el modelo metodológicamente ligado a la reconstrucción.
AR(49) era el benchmark lineal fuerte.
HAR-logRV acaba siendo el modelo operativo más sensato para MVP.